\documentclass[12pt]{article}
\usepackage[utf8]{inputenc}
\usepackage[T1]{fontenc}
\usepackage[margin=1in]{geometry}
\usepackage{setspace}
\usepackage{xcolor}
\usepackage[colorlinks=true, urlcolor=blue!60!black, linkcolor=blue!60!black]{hyperref}
\setlength{\parindent}{0em}
\setlength{\parskip}{0.6em}
\renewcommand{\baselinestretch}{1.0}
\begin{document}
\begin{center}
{\LARGE \textit{United States v. Carolene Products Co.}, 304 U.S. 144 (1938)\par}
\vspace{0.5em}
{\large Case Brief --- Two-Page Edition\par}
\end{center}

\section{Memory Jogger}
Upholding the federal ban on ``filled milk,'' the Court cemented rational-basis review for economic regulation and, in Footnote Four, sketched heightened scrutiny for specific rights, the political process, and ``discrete and insular minorities'' (\href{https://supreme.justia.com/cases/federal/us/304/144/}{304 U.S. 144}).

\section{Detailed Case Facts}
The Filled Milk Act of 1923 barred interstate shipment of skimmed milk blended with non-milk fats to imitate milk, declaring it injurious to health and a fraud on the public. Carolene Products was indicted for shipping ``Milnut'' (condensed skimmed milk and coconut oil) and challenged the Act under the Commerce Clause, the Tenth Amendment, and Fifth Amendment due process.

\section{Procedural History}
The district court sustained a demurrer and dismissed the indictment; the United States appealed directly to the Supreme Court.

\section{Judicial Votes}
Reversed, 6--1. Stone wrote for Hughes, Brandeis, and Roberts; Black joined all but Part Third (so Footnote Four drew four votes); Butler concurred in the result; McReynolds dissented; Cardozo and Reed did not participate.

\section{Holding}
The Act is valid on its face. Congress may exclude articles it reasonably deems harmful from interstate commerce, and economic regulation is presumed constitutional if it rests on any rational basis. Footnote Four is plurality dicta.

\section{Analysis of Opinions}
The commerce power is plenary and does not offend the Tenth Amendment. Due process requires only a rational basis, supported here by the congressional record, and a legislature may strike at one abuse without banning all (here, oleomargarine). Footnote Four suggested ``narrower scope'' for the presumption when laws fall within a specific constitutional prohibition, restrict political processes, or reflect prejudice against discrete and insular minorities. Butler would have let the company prove Milnut wholesome at trial.

\section{Examples: Future Applications}
\textbf{Same-Side:} a degree-and-exam license for interior designers is ordinary economic regulation and survives. \textbf{Opposite-Side:} a permit requirement for leaflets criticizing zoning falls within the First Amendment and burdens the political process, so scrutiny tightens. \textbf{Fence-Sitter:} a ban on a safe, labeled milk substitute backed only by incumbent farmers tests whether any rational basis remains.

\section{Critique}
Libertarians and public-choice scholars see the Act as dairy-lobby protectionism and the presumption as an abdication of judicial duty. Process theorists, led by John Hart Ely, built representation-reinforcement on Footnote Four, though Bruce Ackerman doubted that discrete and insular groups are the least powerful.

\section{Key Quotations}
\begin{quote}
``A legislature may hit at an abuse which it has found, even though it has failed to strike at another.''
\end{quote}
\noindent --- Justice Stone, Opinion of the Court, 304 U.S. at 151.

\end{document}
